\documentclass[11pt]{article}
\usepackage{setspace}
\usepackage{geometry}
\geometry{margin=1in}
\setstretch{1.15}

\begin{document}

% --- Applicant Information ---
\noindent
\textbf{Zhongli Wang} \\
Phone: +49 160 9541 6015 \\
Email: nkwangzhongli@gmail.com \\
Website: https://sites.google.com/view/zhongli-wang \\
Bielefeld, Germany \\

\vspace{0.5cm}

% --- Recipient Information ---
\noindent
\textbf{Prof.\ Mehdi Farsi} \\
Institute of Economic Research \\
University of Neuch\^{a}tel \\
Neuch\^{a}tel, Switzerland \\

\vspace{0.5cm}

% --- Subject Line ---
\noindent
\textbf{Application for Postdoctoral Position in Applied Microeconomics} \\

\vspace{0.3cm}

\noindent Dear Prof.\ Farsi,\\[0.2cm]

I am writing to apply for the postdoctoral position at the Institute of Economic Research, University of Neuch\^{a}tel, starting 1 September 2026. 
I am currently completing my Ph.D.\ in Economics at Bielefeld University and Universit\`a Cattolica del Sacro Cuore (EPOC Marie Sk{\l}odowska--Curie programme). 
My research combines applied microeconomics, behavioral modeling, and advanced computational methods (simulation-based econometrics and machine learning) to study decision-making under frictions and heterogeneity. 
I am particularly interested in extending this work toward energy and environmental economics using household-level evidence and experimental or quasi-experimental designs. 
The opportunity to contribute to projects related to the SWEET consortium and the Swiss household energy survey (SHEDS), while developing a publication pipeline and an independent research profile, is an excellent fit with my goals.

During my Ph.D.\ (started in 2021), I developed structural and simulation-based frameworks to study how agents adjust behavior when incentives, constraints, and information change. 
My dissertation builds quantitative models of technology-intensive markets in which heterogeneous agents learn from experience and peers, update beliefs, and respond to policy and institutional constraints. 
Methodologically, I combine micro data with simulation-based estimation, sensitivity analysis, and out-of-sample validation. 
This background gives me strong preparation for applied work on household energy decisions where adoption, demand responses, and behavioral frictions are central, and where credible identification and careful measurement are essential.

I am excited by the possibility of integrating my modeling skills with SHEDS-style microdata to study household energy demand and the distributional consequences of the energy transition. 
Concretely, I would aim to develop a set of publishable papers that connect behavioral mechanisms (e.g., inattention, present bias, misperceptions about costs/savings, peer effects) to policy-relevant outcomes (e.g., technology adoption, efficiency investments, load shifting, and responses to incentives). 
I am also interested in linking survey measures to revealed behaviors where feasible, and in using experimental designs (survey experiments or field collaborations) to isolate mechanisms that standard observational approaches cannot identify. 
These directions align closely with the position's focus on energy-related applied research and with the goal of building a strong portfolio for the academic job market.

In addition to my dissertation, I have coauthored work that demonstrates my experience with high-dimensional data and reproducible research workflows. 
\textit{Polytope Fraud Theory} applies supervised machine learning to detect structured irregularities in firm financial statements, emphasizing transparent feature design and rigorous evaluation. 
A second project, \textit{Agentic Workflows for Economic Research}, develops a framework for integrating large-language-model agents into empirical pipelines, focusing on documentation, reproducibility, and human-in-the-loop validation. 
Together, these projects reflect my ability to build reliable data pipelines, work with complex datasets, and translate quantitative outputs into clear research claims---skills directly relevant for working with large household surveys and producing policy-relevant research.

I am also motivated by the teaching profile of the position. 
My teaching experience includes two years as a teaching assistant at the University of Zurich, where I supported quantitative business and economics instruction and contributed to the SYLVA project, an automated scoring and feedback system for large courses. 
I would be glad to contribute to the intensive Math refresher for incoming MSc students and to teach the principles course and the master-level Behavioral Economics course. 
My teaching approach emphasizes clarity, structured reasoning, and connecting formal tools to real-world applications---including data-driven exercises in R/Python when appropriate.

Thank you so much for considering my application!
I would be delighted to discuss how my research agenda and quantitative background can contribute to the Institute's work on energy and household behavior, and to the broader research community at the University of Neuch\^{a}tel.

\vspace{6mm}

\noindent Sincerely,\\[12mm]
\noindent Zhongli Wang

\end{document}